\documentclass[12pt]{article}
\usepackage[T1]{fontenc}
\usepackage[utf8]{inputenc}
\usepackage{lmodern}
\usepackage{textcomp}
\usepackage{enumitem}
\usepackage{geometry}
\usepackage{graphicx}
\usepackage{amsmath, amssymb}
\geometry{margin=1in}
\begin{document}
\begin{center}
Constitutional Law - Undergraduate Level Assessment 5 \
Total Marks: 40 \
Answer all questions.
\end{center}

\section*{Multiple Choice (10 marks)}
\begin{enumerate}[label=\arabic*.]
\item What is 'unilateral acts'?
\begin{enumerate}[label=(\alph*)]
    \item They are acts that States perform as practice in the context of custom
    \item They are acts creating unilateral legal obligations to the acting State
    \item Unilateral acts are simply political acts of State devoid of any legal effect
    \item Unilateral acts are those that State perform in order to be bound by a treaty
\end{enumerate}
\item Why does communitarianism resemble Aristotle's philosophy more than Kant's?
\begin{enumerate}[label=(\alph*)]
    \item Because Aristotle justified slavery.
    \item Because Kant failed to distinguish individual from social morality?
    \item Because Aristotle believed that man is a 'social animal'
    \item Because Kant regarded the individual as unimportant.
\end{enumerate}
\item What is Kant's 'categorical imperative?
\begin{enumerate}[label=(\alph*)]
    \item You must not consciously harm another person.
    \item You must always act in the best interests of the community.
    \item You must treat human beings as means rather than ends.
    \item You must act according as if your values apply to everyone.
\end{enumerate}
\item Maine's famous aphorism that 'the movement of progressive societies has hitherto been a movement from Status to Contract' is often misunderstood. In what way?
\begin{enumerate}[label=(\alph*)]
    \item It is misinterpreted as a prediction.
    \item His concept of status is misrepresented.
    \item It is taken literally.
    \item His idea is considered inapplicable to Western legal systems.
\end{enumerate}
\item Which is the least accurate description of legal positivism?
\begin{enumerate}[label=(\alph*)]
    \item It regards morals and law as inseparable.
    \item It perceives law as commands.
    \item It regards a legal order as a closed logical system.
    \item It espouses the view that there is no necessary connection between morality and law.
\end{enumerate}
\end{enumerate}

\section*{True / False (10 marks)}
\textit{Answer True or False}
\begin{enumerate}[label=\arabic*.]
\setcounter{enumi}{5}
\item True or False: Austin argues that the science of jurisprudence focuses mainly on morality and ethical principles.
\item True or False: The Universal Declaration of Human Rights (UDHR) is a resolution adopted by the UN General Assembly.
\item True or False: The protective principle of jurisdiction is based on the harm to national interests caused by conduct committed abroad.
\item True or False: Adherence to precedent in legal cases is known as the doctrine of stare decisis.
\item True or False: Individuals possess unlimited international legal personality, allowing them to act independently on the global stage.
\end{enumerate}

\section*{Long Form Response (20 marks)}
\textit{Answer each question with a clear and complete explanation.}
\begin{enumerate}[label=\arabic*.]
\setcounter{enumi}{10}
\item Discuss the Social Fact Thesis in positivism as articulated by Ken Himma, and analyze its implications for understanding law as a social construction.
\item Critically analyze the 'fair play' argument for obeying the law and discuss the implications of an unfair legal system on this duty.
\end{enumerate}

\vfill
\noindent\textit{End of Paper}
\end{document}